\section{Contribution}\label{section::contribution}

A primary contribution of this thesis is the development and public release of a large-scale multilingual dataset specifically structured for the task of text-to-SPARQL generation.
Comprising approximately 895,954 examples, this dataset spans 12 languages and is derived from established knowledge extraction benchmarks, including the known QALD series \cite{usbeck2024qald}, alongside other relevant challenges.
The creation process involved meticulous data compilation from these diverse sources, followed by cleaning to remove inconsistencies and noise often present in aggregated datasets.
Furthermore, the data was standardized into a consistent question-SPARQL query format and structured to facilitate efficient supervised fine-tuning of Large Language Models.
A significant aspect of this dataset is its augmentation, particularly through the generation of German translations for existing English questions and the integration of Wikidata entity and relationship mappings.
This augmentation not only expands the linguistic coverage but also provides explicit grounding crucial for the entity linking sub-task, a known challenge in text-to-SPARQL systems \cite{brei2024leveraging}.
To promote reproducibility and further research within the community, the complete dataset and the scripts used for its generation, have been made publicly available on the Hugging Face Hub \cite{lhoest2021datasets} and Github respectively.

\todo{scripts to replicate the whole data extraction, transform, loading process, partly automated given a repo by scraping}

\todo{8 finetuned models, one for each language and training approach.}

\todo{different training approaches like using bleu metric as evaluation metric, stop loss and so on}

\todo{several mistakes made along the way such as overall training with the whole dataset}

\todo{results for 3 different benchmarks}

\todo{insights into why we get what we are getting}

\todo{suggestions for several venues of probable research pathways}

\todo{pitfalls as well}